\subsection{外乱抑制とフィードフォワード}

外乱または目標値の変化が測定可能である場合、フィードフォワードを加えることで偏差伝達を低減できる。対象が $Y=G U+G_dD$ と書け、外乱 $D$ が測定可能であるなら
\begin{equation}
  U=C(R-Y)+F_dD
\end{equation}
とし、理想的には
\begin{equation}
  F_d(s)=-G(s)^{-1}G_d(s)
\end{equation}
を選ぶ。ただし $G^{-1}$ が不安定、非プロパー、または不確かな場合は直接実装できない。その場合には、低周波成分のみを補償する近似や入力制限を含む定式化を用いる。

目標値フィードフォワードも同じ考え方である。位置制御で目標軌道 $r(t)$ の速度や加速度が分かっていれば、必要な力を
\begin{equation}
  u_{ff}=m\ddot{r}+c\dot{r}+kr
\end{equation}
として与え、フィードバックでモデル誤差と外乱に起因する偏差を補償する。

\subsection{周波数領域を支える定理}

以下では、正弦波定常応答において伝達関数へ $s=j\omega$ を代入する根拠を示す。

\begin{theorem}[正弦波定常応答の定理]
安定な線形時不変系の伝達関数を $G(s)$ とする。入力 $u(t)=\operatorname{Re}\{Ue^{j\omega t}\}$ を加えると、過渡応答が消えた後の出力は
\begin{equation}
  y_{ss}(t)=\operatorname{Re}\{G(j\omega)Ue^{j\omega t}\}
\end{equation}
となる。
\end{theorem}

\begin{derivation}
複素指数入力 $u_c(t)=Ue^{j\omega t}$ を考える。線形時不変系では微分作用素 $\dd/\dd t$ に対して
\begin{equation}
  \frac{\dd}{\dd t}e^{j\omega t}=j\omega e^{j\omega t}
\end{equation}
である。たとえば微分方程式
\begin{equation}
  a_n y^{(n)}+\cdots+a_1\dot{y}+a_0y
  =b_m u^{(m)}+\cdots+b_0u
\end{equation}
に $y_c(t)=Ye^{j\omega t}$、$u_c(t)=Ue^{j\omega t}$ を代入する。$k$ 階微分は $(j\omega)^k$ 倍になるので
\begin{equation}
  \{a_n(j\omega)^n+\cdots+a_0\}Ye^{j\omega t}
  =\{b_m(j\omega)^m+
  \cdots+b_0\}Ue^{j\omega t}
\end{equation}
である。両辺を $e^{j\omega t}$ で割ると
\begin{equation}
  \frac{Y}{U}=\frac{b_m(j\omega)^m+\cdots+b_0}{a_n(j\omega)^n+\cdots+a_0}=G(j\omega)
\end{equation}
となる。したがって $Y=G(j\omega)U$ であり、実際の正弦波応答はその実部である。
\end{derivation}

\begin{formula}[ボード線図の積和変換]
複数の因子の積 $G(s)=G_1(s)G_2(s)$ に対して
\begin{equation}
  20\log_{10}\abs{G(j\omega)}
  =20\log_{10}\abs{G_1(j\omega)}+20\log_{10}\abs{G_2(j\omega)}
\end{equation}
である。また位相は
\begin{equation}
  \angle G(j\omega)=\angle G_1(j\omega)+\angle G_2(j\omega)
\end{equation}
である。
\end{formula}

\begin{derivation}
絶対値の性質 $\abs{ab}=\abs{a}\abs{b}$ と対数の性質 $\log(ab)=\log a+\log b$ より
\begin{align}
  20\log_{10}\abs{G_1G_2}
  &=20\log_{10}\{\abs{G_1}\abs{G_2}\}\\
  &=20\log_{10}\abs{G_1}+20\log_{10}\abs{G_2}.
\end{align}
複素数を $G_i=\abs{G_i}e^{j\phi_i}$ と書けば
\begin{equation}
  G_1G_2=\abs{G_1}\abs{G_2}e^{j(\phi_1+\phi_2)}
\end{equation}
なので、位相も和になる。
\end{derivation}

\begin{theorem}[ナイキスト安定判別]
単位負帰還系の開ループ $L(s)$ が虚軸上に極を持たず、$1+L(j\omega)\ne0$ であるとする。右半平面を時計回りに囲むナイキスト経路を用い、開ループ $L(s)$ の右半平面極の数を $P$、ナイキスト線図が点 $-1$ を時計回りに囲む回数を $N$ と数える。このとき閉ループ特性方程式 $1+L(s)=0$ の右半平面根の数 $Z$ は
\begin{equation}
  Z=N+P
\end{equation}
である。閉ループ安定性には $Z=0$ が必要十分である。
\end{theorem}

この定理は複素解析の偏角原理に基づく。反時計回りの囲みを正に数える流儀では符号が変わるため、ナイキスト判別を使うときは $N$ の符号規約を必ず明示する。ナイキスト線図は、閉ループ極の個数を周波数応答から決定するために用いられる。

\begin{formula}[根軌跡の角度条件と大きさ条件]
閉ループ特性方程式を $1+KG(s)=0$ とする。$K>0$ の根軌跡上の点 $s$ は
\begin{align}
  \angle G(s)&=(2\ell+1)\pi,\qquad \ell\in\mathbb{Z},\\
  K&=\frac{1}{\abs{G(s)}}
\end{align}
を満たす。
\end{formula}

\begin{derivation}
$1+KG(s)=0$ より
\begin{equation}
  KG(s)=-1
\end{equation}
である。右辺 $-1$ は大きさが 1、偏角が $(2\ell+1)\pi$ の複素数である。$K$ は正の実数なので偏角を変えない。したがって
\begin{equation}
  \angle G(s)=(2\ell+1)\pi
\end{equation}
でなければならない。また大きさを比べると
\begin{equation}
  K\abs{G(s)}=1
\end{equation}
だから
\begin{equation}
  K=1/\abs{G(s)}
\end{equation}
を得る。
\end{derivation}
